\chapter{Construction du système de scoring expert}

\section{Principe de la grille}

Le score expert vise à formaliser, sous forme d'une grille pondérée, un jugement de type
« analyste de crédit » à partir de ratios financiers directement calculables sur les variables
canoniques disponibles (Chapitre~3). Pour chaque ratio retenu, une valeur normalisée sur
$[0, 100]$ est calculée selon deux bornes de saturation (plancher/plafond) et un sens
d'interprétation (« plus est mieux » ou « moins est mieux »), puis agrégée par une moyenne
pondérée.

\begin{definition}{Score expert}
Soit $R = \{r_1, \dots, r_k\}$ l'ensemble des ratios financiers retenus, $w_i$ le poids associé
à $r_i$, et $s_i \in [0, 100]$ le score normalisé du ratio $r_i$ pour un dossier donné. Le score
expert est défini par :
\[
S_{\text{expert}} = \frac{\sum_{i \,:\, r_i \text{ disponible}} w_i \cdot s_i}{\sum_{i \,:\, r_i \text{ disponible}} w_i}
\]
Le score n'est calculé que sur les ratios effectivement disponibles pour le dossier (pondération
renormalisée), ce qui permet de traiter les dossiers incomplets sans imputation arbitraire.
\end{definition}

\section{Ratios retenus et pondération}

Six ratios ont été retenus (Tableau~\ref{tab:grille-experte}), couvrant les dimensions de
rentabilité, de solvabilité et de liquidité — trois piliers classiques de l'analyse financière
du risque de crédit (cf. cadre théorique, Chapitre~2).

\begin{table}[h]
\centering
\small
\begin{tabular}{@{}lccc@{}}
\toprule
\textbf{Ratio} & \textbf{Poids} & \textbf{Sens} & \textbf{Bornes} \\
\midrule
Rentabilité financière & 20\,\% & plus = mieux & $[-20, 40]$ \\
Autonomie financière & 20\,\% & plus = mieux & $[0, 100]$ \\
Endettement à court terme & 15\,\% & moins = mieux & $[0, 100]$ \\
Liquidité réduite & 15\,\% & plus = mieux & $[0, 200]$ \\
Capacité de remboursement & 15\,\% & plus = mieux & $[0, 10]$ \\
Dettes nettes / fonds propres & 15\,\% & moins = mieux & $[0, 5]$ \\
\bottomrule
\end{tabular}
\caption{Pondération de la grille experte (hypothèse de travail)}
\label{tab:grille-experte}
\end{table}

\begin{limite}[Pondérations non calibrées]
Ces pondérations constituent une \emph{hypothèse de travail initiale}, construite par analogie
avec les grilles de scoring classiques documentées dans la littérature (Chapitre~2), et non une
reproduction de la grille interne réelle de la Banque Populaire, qui n'est pas documentée dans
le périmètre accessible pour ce stage. Une calibration avec le maître de stage — par exemple par
ajustement des poids jusqu'à minimiser l'écart avec la notation MNS2 réelle sur les dossiers
disponibles — constitue une piste d'amélioration immédiate (cf. Chapitre~7).
\end{limite}

\section{Résultats}

Le Tableau~\ref{tab:resultats-grille} présente le score expert calculé pour les neuf dossiers de
l'échantillon.

\begin{table}[h]
\centering
\small
\begin{tabular}{@{}lc@{}}
\toprule
\textbf{Dossier} & \textbf{Score expert / 100} \\
\midrule
Feuil13 & 100.0 \\
Feuil2  & 76.4 \\
Feuil5  & 60.5 \\
Feuil8  & 53.8 \\
Feuil12 & 51.5 \\
Feuil1  & 44.2 \\
Feuil11 & 42.3 \\
Feuil3  & 39.2 \\
Feuil4  & N/A (données insuffisantes) \\
\bottomrule
\end{tabular}
\caption{Score expert par dossier}
\label{tab:resultats-grille}
\end{table}

Le score de 100/100 obtenu par Feuil13 mérite un commentaire : il résulte du fait que ce dossier
ne dispose que d'un très petit nombre de ratios renseignés parmi les six retenus, et que ceux-ci
sont favorables. Ce cas illustre concrètement une limite du mécanisme de renormalisation évoqué
en Section~4.1 : un score calculé sur peu de ratios disponibles est statistiquement moins fiable
qu'un score calculé sur l'ensemble de la grille, sans que cela soit visible dans le chiffre final.

\begin{limite}
Le score expert ne porte, à ce stade, aucune indication de \emph{confiance} liée au nombre de
ratios effectivement disponibles. Une évolution naturelle consisterait à afficher, à côté du
score, le nombre de ratios mobilisés (ex. « 51,5/100, calculé sur 4/6 ratios »).
\end{limite}

\section{Synthèse}

Le score expert fournit une base transparente et immédiatement justifiable, conforme aux
exigences opérationnelles identifiées au Chapitre~2 (Section 2.2). Le Chapitre~5 construit, à
partir des mêmes données, un score par apprentissage automatique, avant que le Chapitre~6 ne
confronte les deux approches à la notation bancaire réelle.
